\chapter{Conclusiones tecnicas}
\label{chap:conclusiones-tecnicas}

\section{Estado tecnico del proyecto}
\label{sec:conclusiones-estado}

La plataforma Robot Explota Globos cuenta con una arquitectura Django funcional para gestionar informacion publica del torneo, registro de participantes, confirmacion de asistencia, sincronizacion a equipos, control operativo por division, fases de competencia, resultados, correcciones manuales, historial y comunicaciones.

El proyecto no es solo una coleccion de vistas. El dominio central esta modelado alrededor de ediciones, registros, equipos, competencias por division, grupos, batallas, estados e historial. Esa estructura permite comprender que el sistema resuelve un problema operativo especifico: mantener trazabilidad y control de una competencia que puede cambiar durante el evento sin perder consistencia de datos.

\section{Arquitectura alcanzada}
\label{sec:conclusiones-arquitectura}

La arquitectura alcanzada separa responsabilidades de forma razonable:

\begin{itemize}
  \item \archivo{apps/core} concentra paginas publicas;
  \item \archivo{apps/participants} administra registros, asistencia y sincronizacion a equipos;
  \item \archivo{apps/tournament} concentra modelo competitivo, servicios, panel y reglas de avance;
  \item \archivo{apps/invitations} gestiona comunicaciones, plantillas, destinatarios, lotes y logs;
  \item \archivo{config} centraliza configuracion por entorno;
  \item templates y JavaScript actuan como capa de operacion, no como autoridad del dominio.
\end{itemize}

La separacion entre vistas y servicios es especialmente relevante. Las vistas reciben solicitudes y preparan respuestas; los servicios contienen reglas de torneo, persistencia de estados, correcciones e historial. Esto mejora la mantenibilidad porque permite razonar sobre el flujo competitivo sin perseguir toda la logica dentro de templates o JavaScript.

\section{Mantenibilidad}
\label{sec:conclusiones-mantenibilidad}

El sistema es mantenible si se respetan sus fronteras:

\begin{itemize}
  \item cambios en reglas de torneo deben revisarse en servicios y pruebas;
  \item cambios en templates internos deben revisarse junto con \archivo{control.js};
  \item cambios de configuracion deben validar seguridad y entorno;
  \item cambios de comunicaciones deben probar dry-run, test y envio oficial;
  \item cambios de modelos deben acompanar migraciones, respaldo y validacion de datos.
\end{itemize}

La principal deuda tecnica es la concentracion de logica en \archivo{apps/tournament/services.py}. No impide operar el sistema, pero aumenta costo de modificacion. La segunda deuda es la cobertura de pruebas: el manual documenta controles y flujos, pero la evolucion futura exige ampliar pruebas automaticas.

\section{Capacidad de evolucion}
\label{sec:conclusiones-evolucion}

La plataforma puede evolucionar hacia nuevas divisiones, formatos, vistas publicas de resultados, exportaciones e integraciones. Esa capacidad se fundamenta en modelos persistentes, competencia por division, historial, configuracion flexible y servicios de materializacion de fases.

La evolucion debe ser incremental. Agregar formatos sin dividir servicios ni aumentar pruebas puede degradar la estabilidad. Publicar resultados en tiempo real sin resolver despliegue, seguridad y monitoreo puede exponer informacion o estados parciales. Operar comunicaciones reales sin politica de datos y backups aumenta riesgo institucional.

\section{Condiciones para operacion real}
\label{sec:conclusiones-operacion}

Antes de operar con datos reales y participantes reales, deben cerrarse estas condiciones:

\begin{itemize}
  \item definir infraestructura de produccion;
  \item configurar dominio, HTTPS, hosts y CSRF;
  \item usar PostgreSQL con backups;
  \item probar SMTP y LibreOffice;
  \item ejecutar simulacion completa del torneo;
  \item definir responsables y permisos operativos;
  \item formalizar tratamiento de datos personales;
  \item ampliar pruebas para flujos criticos.
\end{itemize}

Estas condiciones no son decorativas. Cada una protege una parte concreta del sistema: configuracion, datos, comunicaciones, competencia, seguridad y recuperacion.

\section{Beneficios del diseno}
\label{sec:conclusiones-beneficios}

El diseno aporta beneficios concretos:

\begin{itemize}
  \item reduce duplicacion al centralizar reglas en servicios;
  \item conserva historial de movimientos y resultados;
  \item permite correcciones manuales controladas;
  \item separa informacion publica de operacion interna;
  \item admite simulacion y prueba controlada de comunicaciones;
  \item diferencia desarrollo local de produccion mediante variables;
  \item mantiene una base documental suficiente para transferencia tecnica.
\end{itemize}

El resultado es una plataforma con base tecnica solida para el contexto del torneo, siempre que las observaciones de pruebas, despliegue, backups y privacidad se atiendan antes de produccion.

\section{Recomendaciones para futuras versiones}
\label{sec:conclusiones-recomendaciones}

\begin{itemize}
  \item Priorizar pruebas integrales del flujo de competencia y permisos internos.
  \item Dividir servicios de torneo por responsabilidad antes de ampliar formatos.
  \item Cerrar guia de despliegue real cuando se confirme infraestructura.
  \item Formalizar backups, privacidad, responsables y calendario operativo.
  \item Mantener el manual como documento vivo: cada cambio en modelo, servicio, template o contrato AJAX debe reflejarse en la documentacion.
\end{itemize}
